\documentclass[tikz,border=4pt]{standalone}
\usepackage{ctex}
\usepackage{amsmath,amssymb}
\usepackage{pgfplots}
\pgfplotsset{compat=1.18}
\usetikzlibrary{decorations.markings}
\begin{document}
\begin{tikzpicture}
% ====== 左：第一类（对弧长）方向无关 ======
\begin{axis}[
  name=left,
  width=7cm, height=6cm,
  xmin=-0.5, xmax=4.5, ymin=-2.0, ymax=3,
  xlabel={$x$}, ylabel={$y$},
  axis lines=center,
  xlabel style={right}, ylabel style={above},
  title={\textbf{第一类（对弧长 $ds$）}},
  title style={font=\small},
  clip=false,
]
% L 正向
\addplot[blue!70!black, very thick, samples=80, domain=0.3:3.7,
  postaction={decorate, decoration={markings,
    mark=at position 0.55 with {\arrow{>}}}}]
  {0.5 + 1.2*sin(deg(x*0.9))};
\node[blue!70!black, font=\small] at (axis cs: 3.9, 1.6) {$L$};

% L^- 反向，下移虚线
\addplot[blue!50, very thick, dashed, samples=80, domain=3.7:0.3,
  postaction={decorate, decoration={markings,
    mark=at position 0.55 with {\arrow{>}}}}]
  {0.0 + 1.2*sin(deg(x*0.9))};
\node[blue!50!black, font=\footnotesize, anchor=west]
  at (axis cs: 0.0, -1.3) {$L^-$ (反向)};

% ds 标记
\node[red, font=\small] at (axis cs: 2.0, 1.55) {$\bullet$};
\draw[red, thick] (axis cs: 2.0, 1.45) -- (axis cs: 2.4, 1.6);
\node[red, font=\footnotesize, anchor=south west] at (axis cs: 2.4, 1.6) {$ds>0$};

% 结论框
\node[draw=gray!50, fill=gray!10, fill opacity=0.85,
      text opacity=1, rounded corners=2pt, font=\footnotesize, align=center]
  at (axis cs: 2.1, -1.85)
  {$\displaystyle\int_{L^-}\!f\,ds=\int_L f\,ds$\\
   \textbf{方向无关}（$ds\ge 0$）};
\end{axis}

% ====== 右：第二类（对坐标）方向相关 ======
\begin{axis}[
  at={(left.east)}, xshift=2.2cm, anchor=west,
  width=7cm, height=6cm,
  xmin=-0.5, xmax=4.5, ymin=-2.0, ymax=3,
  xlabel={$x$}, ylabel={$y$},
  axis lines=center,
  xlabel style={right}, ylabel style={above},
  title={\textbf{第二类（对坐标 $dx,dy$）}},
  title style={font=\small},
  clip=false,
]
% 曲线 + 方向箭头
\addplot[orange!80!black, very thick, samples=80, domain=0.3:3.7,
  postaction={decorate, decoration={markings,
    mark=at position 0.55 with {\arrow{>}}}}]
  {0.5 + 1.2*sin(deg(x*0.9))};
\node[orange!80!black, font=\small] at (axis cs: 3.9, 1.6) {$L$};

% 切向量 T 在三点（预计算）
% t=0.8: y=0.5+1.2*sin(41.25°)=1.291, slope=1.08*cos(41.25°)=0.812, norm=1.290
%        dx=0.426, dy=0.346 → 终点 (1.226, 1.637)
\draw[->, very thick, red] (axis cs: 0.8, 1.291) -- (axis cs: 1.226, 1.637);
% t=1.8: y=0.5+1.2*sin(92.81°)=1.699, slope=1.08*cos(92.81°)=-0.053, norm=1.001
%        dx=0.549, dy=-0.029 → 终点 (2.349, 1.670)
\draw[->, very thick, red] (axis cs: 1.8, 1.699) -- (axis cs: 2.349, 1.670);
% t=3.0: y=0.5+1.2*sin(154.7°)=1.012, slope=1.08*cos(154.7°)=-0.977, norm=1.398
%        dx=0.393, dy=-0.384 → 终点 (3.393, 0.628)
\draw[->, very thick, red] (axis cs: 3.0, 1.012) -- (axis cs: 3.393, 0.628);
\node[red, font=\footnotesize, anchor=north east]
  at (axis cs: 4.0, 0.3) {$\mathbf{T}$ (切向)};

% 结论框
\node[draw=gray!50, fill=gray!10, fill opacity=0.85,
      text opacity=1, rounded corners=2pt, font=\footnotesize, align=center]
  at (axis cs: 2.1, -1.85)
  {$\displaystyle\int_{L^-}\!P\,dx+Q\,dy=-\!\int_L P\,dx+Q\,dy$\\
   \textbf{方向相关}（$dx$ 带符号）};
\end{axis}
\end{tikzpicture}
\end{document}
